% =====================================================================
%  config/comandos.tex
%  Chaves do ciclo, caixas de apoio e macros do relatório.
%  A equipe não precisa mexer aqui.
% =====================================================================

% ---------------------------------------------------------------------
%  1. CICLO DE PROJETO (AP1, AP2 ou AP3)
% ---------------------------------------------------------------------
\ifdefined\CICLO\else\def\CICLO{1}\fi
\newcounter{ciclo}
\setcounter{ciclo}{\CICLO}

% \seAP{n}{texto}     -> aparece somente no ciclo n
% \aPartirDaAP{n}{t}  -> aparece do ciclo n em diante
\newcommand{\seAP}[2]{\ifnum\value{ciclo}=#1\relax#2\fi}
\newcommand{\aPartirDaAP}[2]{\ifnum\value{ciclo}<#1\relax\else#2\fi}

\newcommand{\nomeciclo}{%
  \ifcase\value{ciclo}\or
    Ciclo de Projeto 1 --- sistema embarcado de sensoriamento e atuação\or
    Ciclo de Projeto 2 --- integração à infraestrutura de comunicação e
      formação do conjunto de dados\or
    Ciclo de Projeto 3 --- projeto integrador de AIoT\fi}

\newcommand{\rotuloAP}{AP\arabic{ciclo}}

\newcommand{\moduloDoCiclo}{%
  \ifcase\value{ciclo}\or Módulo 1\or Módulo 2\or Módulo 3\fi}

% ---------------------------------------------------------------------
%  2. NOTAS AO AUTOR (instruções que somem na versão final)
% ---------------------------------------------------------------------
\newif\ifmostrarnotas
\ifdefined\MOSTRARNOTAS\mostrarnotastrue\else\mostrarnotasfalse\fi

\newtcolorbox{tcbnota}{
  enhanced, breakable,
  colback   = gray!6,
  colframe  = gray!45,
  boxrule   = 0.4pt,
  left      = 6pt, right = 6pt, top = 5pt, bottom = 5pt,
  fontupper = \small\sffamily,
  borderline west = {2pt}{0pt}{gray!55},
  sharp corners,
  before skip = 8pt, after skip = 8pt,
}

\usepackage{comment}
\ifmostrarnotas
  \newenvironment{orientacao}{\begin{tcbnota}\noindent\textbf{Como preencher.}\ }%
                       {\end{tcbnota}}
  \newenvironment{exemploguia}{\begin{tcbnota}\noindent\textbf{Exemplo.}\ }%
                          {\end{tcbnota}}
\else
  \excludecomment{orientacao}
  \excludecomment{exemploguia}
\fi

% ---------------------------------------------------------------------
%  3. DECISÃO DE PROJETO
%     Vale 25% da nota da AP. Uma caixa por decisão que a equipe tomou.
% ---------------------------------------------------------------------
\newcounter{decisao}
\renewcommand{\thedecisao}{DP-\arabic{decisao}}

\newtcolorbox[use counter=decisao]{caixadecisao}[1]{
  enhanced, breakable,
  colback    = corDestaque!4,
  colframe   = corDestaque!70,
  boxrule    = 0.6pt,
  left       = 7pt, right = 7pt, top = 6pt, bottom = 6pt,
  sharp corners,
  fonttitle  = \bfseries\sffamily\small,
  title      = {DP-\thetcbcounter\ --- #1},
  before skip = 10pt, after skip = 10pt,
}

\newenvironment{decisao}[1]{\begin{caixadecisao}{#1}\small}{\end{caixadecisao}}

\newcommand{\campoDecisao}[2]{\par\smallskip\noindent\textbf{#1}\ #2}
\newcommand{\alternativas}[1]{\campoDecisao{Alternativas consideradas.}{#1}}
\newcommand{\criterio}[1]{\campoDecisao{Critério de escolha.}{#1}}
\newcommand{\escolha}[1]{\campoDecisao{Decisão.}{#1}}
\newcommand{\evidencia}[1]{\campoDecisao{Evidência (número e procedência).}{#1}}
\newcommand{\consequencia}[1]{\campoDecisao{Consequência assumida.}{#1}}

% ---------------------------------------------------------------------
%  4. IDENTIFICADORES E VEREDITOS
% ---------------------------------------------------------------------
\newcommand{\RF}[1]{\mbox{\textbf{RF-#1}}}      % requisito funcional
\newcommand{\RNF}[1]{\mbox{\textbf{RNF-#1}}}    % requisito não funcional
\newcommand{\CT}[1]{\mbox{\textbf{CT-#1}}}      % caso de teste
\newcommand{\DP}[1]{\mbox{\textbf{DP-#1}}}      % decisão de projeto

\definecolor{corAprovado}{HTML}{1B6B3A}
\definecolor{corReprovado}{HTML}{9B1C1C}
\definecolor{corParcial}{HTML}{8A6100}
\newcommand{\aprovado}{\textcolor{corAprovado}{\textbf{atende}}}
\newcommand{\reprovado}{\textcolor{corReprovado}{\textbf{não atende}}}
\newcommand{\parcial}{\textcolor{corParcial}{\textbf{atende em parte}}}

% ---------------------------------------------------------------------
%  5. ETIQUETA DA RUBRICA
%     Marca discretamente a que fatia da nota a seção responde.
% ---------------------------------------------------------------------
\newcommand{\rubricaTag}[2]{%
  \par\vspace{-6pt}\noindent
  {\footnotesize\sffamily\color{gray!75}%
   Rubrica da \rotuloAP: #1 --- #2\,\%}%
  \par\vspace{4pt}}

\newcommand{\rubricaFuncionalidade}{\rubricaTag{funcionalidade do protótipo}{30}}
\newcommand{\rubricaJustificativa}{\rubricaTag{justificativa técnica das decisões}{25}}
\newcommand{\rubricaRobustez}{\rubricaTag{robustez, testes e documentação}{25}}
\newcommand{\rubricaComunicacao}{\rubricaTag{comunicação dos resultados}{20}}

% ---------------------------------------------------------------------
%  6. ATALHOS DE TABELA
% ---------------------------------------------------------------------
\newcolumntype{L}[1]{>{\raggedright\arraybackslash}p{#1}}
\newcolumntype{C}[1]{>{\centering\arraybackslash}p{#1}}
\newcolumntype{R}[1]{>{\raggedleft\arraybackslash}p{#1}}
\newcommand{\cabecalholinha}{\rowcolor{gray!12}}
\renewcommand{\arraystretch}{1.25}

% Procedência de um número: folha de dados, bancada ou estimativa.
\newcommand{\estimativa}{\textcolor{corParcial}{\textsc{\small estimativa}}}
\newcommand{\bancada}{\textcolor{corAprovado}{\textsc{\small bancada}}}
\newcommand{\folhadedados}{\textsc{\small folha de dados}}

% ---------------------------------------------------------------------
%  7. MARCAS DA MATRIZ DE DECISÃO (AP2)
%     Três marcas apenas: os critérios não se convertem em pontos.
% ---------------------------------------------------------------------
\newcommand{\serve}{\textcolor{corAprovado}{\textbf{serve}}}
\newcommand{\ressalva}{\textcolor{corParcial}{\textbf{com ressalva}}}
\newcommand{\naoserve}{\textcolor{corReprovado}{\textbf{não serve}}}

% ---------------------------------------------------------------------
%  8. PENDÊNCIAS DA EQUIPE (rascunho da AP2)
%     Marca em vermelho o que depende de um dado ou decisão da equipe.
%     Antes de entregar: grep -rn "confirmar{" conteudo apendices
%     tem de voltar vazio.
% ---------------------------------------------------------------------
\newcommand{\confirmar}[1]{\textcolor{corReprovado}{\textbf{[confirmar:}~#1\textbf{]}}}
